\begin{thebibliography}{99}

\bibitem{takum}
L.~Hunhold.
\newblock Beating posits at their own game: takum arithmetic.
\newblock arXiv:2404.18603.

\bibitem{tekum}
L.~Hunhold.
\newblock Tekum: balanced-ternary tapered-precision arithmetic.
\newblock arXiv:2512.10964.

\bibitem{takumrtl}
L.~Hunhold.
\newblock Design and implementation of a takum arithmetic hardware codec.
\newblock arXiv:2408.10594. Source: \url{https://github.com/takum-arithmetic/Takum-Codec-RTL}.

\bibitem{catalogue}
Author names omitted for double-blind review.
\newblock An 83-format numeric catalog with bit-exact conformance vectors: a
vendor-neutral reference for FP8, BF16, MXFP4, and microscaling formats.
\newblock arXiv:2606.09686, 2026; v1 8 June 2026, v2 22 June 2026.
\newblock DOI \url{https://doi.org/10.48550/arXiv.2606.09686}.
\newblock The shipped catalog contains 83 formats, matching the title of the
current version. The tagged source release accompanying that paper is public.

\bibitem{etiemble2019} D.~Etiemble, ``Ternary circuits: why R=3 is not the
Optimal Radix for Computation,'' arXiv:1908.06841, 2019.

\bibitem{archhdl2026} S.~Zhao, ``Formally Verified Synthesizable Floating-Point
Data Types in ARCH HDL,'' arXiv:2607.23715, 2026.

\bibitem{mohanty2025} H.~Mohanty, V.~N.~Viswambharan, and D.~N.~Gadde, ``Formal
that ``Floats'' High: Formal Verification of Floating Point Arithmetic,''
arXiv:2512.06850, 2025.

\bibitem{fasoli2026} A.~Fasoli, M.~Kar, C.-C.~Liu, S.~Venkataramani,
V.~Srinivasan, L.~Chang, and N.~Wang, ``Is Finer Better? The Limits of
Microscaling Formats in Large Language Models,'' arXiv:2601.19026, 2026.

\bibitem{cankaya2026} N.~Cankaya, ``Bit-Exact AI Inference Verification Without
Performance Tradeoffs,'' arXiv:2606.00279, 2026.

\bibitem{gond2026llm42} R.~Gond, A.~K.~Kamath, R.~Ramjee, and A.~Panwar,
``LLM-42: Enabling Determinism in LLM Inference with Verified Speculation,''
arXiv:2601.17768, 2026.

\bibitem{schoenbaum2021} L.~T.~Schoenbaum, ``A Tapered Floating Point Extension
for the Redundant Signed Radix 2 System Using the Canonical Recoding,''
arXiv:2105.14401, 2021.

\bibitem{hunhold2024sparse} L.~Hunhold and J.~Quinlan, ``Evaluation of
Bfloat16, Posit, and Takum Arithmetics in Sparse Linear Solvers,''
arXiv:2412.20268, 2024.

\bibitem{mlformats} Performance-Efficiency Trade-off of Low-Precision Numerical Formats in Deep Neural Networks, arXiv:1903.10584.

\bibitem{qiao2026curvefp} Y.~Qiao, ``CurveFP: Co-Designing Numerical Representation and Product Arithmetic for Language Models,'' arXiv:2608.10010, 8~August 2026. \url{https://arxiv.org/abs/2608.10010}

\bibitem{ieee754_2019} IEEE, ``IEEE Standard for Floating-Point Arithmetic,''
IEEE Std 754-2019 (revision of IEEE Std 754-2008), July 2019,
doi:10.1109/IEEESTD.2019.8766229.

\bibitem{p3109wg2026} IEEE Working Group P3109, ``Interim Report on Binary
Floating-point Formats for Machine Learning,'' Version~4.0, 26 June 2026.
\url{https://github.com/P3109/Public}

\bibitem{flops2026} T.-C.~Chang, S.~Park, J.~P.~Lim, and S.~Nagarakatte,
``FLoPS: Semantics, Operations, and Properties of P3109 Floating-Point
Representations in Lean,'' arXiv:2602.15965, 2026.

\bibitem{wintersteiger2025} C.~M.~Wintersteiger, ``Formal Verification of the
IEEE P3109 Standard for Binary Floating-Point Formats for Machine Learning,'' in
\emph{2025 IEEE 32nd Symp.\ Computer Arithmetic (ARITH)}, El Paso, TX, 2025,
pp.~157--166, doi:10.1109/ARITH64983.2025.00032.

\bibitem{fitzgibbon2026kappa} A.~Fitzgibbon, C.~M.~Wintersteiger, and
J.~Sarnoff, ``Novel Aspects of IEEE SA P3109 Arithmetic Formats for Machine
Learning,'' arXiv:2606.04028, 2026.

\bibitem{matsuoka2026fp8a} S.~Matsuoka, ``FP8 is All You Need (Part~1):
Debunking Hardware FP64 as the HPC Holy Grail,'' arXiv:2606.06510, 2026.

\bibitem{matsuoka2026fp8b} S.~Matsuoka, ``FP8 is All You Need (Part~2):
Efficient Ozaki--Bailey Style FFT Through Tensor-core Garner Reformulation and
Kulisch Escape Route,'' arXiv:2606.23698, 2026.

\bibitem{int_vs_fp2025} M.~Chen, M.~Wu, H.~Jin, Z.~Yuan, et al., ``INT vs.\ FP:
A Comprehensive Study of Fine-Grained Low-bit Quantization Formats,''
arXiv:2510.25602, 2025.

\bibitem{tereffic2025} C.~Yin, Z.~Bai, P.~Venkatram, S.~Aggarwal, et al.,
``TerEffic: Highly Efficient Ternary LLM Inference on FPGA,'' arXiv:2502.16473,
2025.

\bibitem{tenet2025} Z.~Huang, R.~Ma, S.~Cao, R.~Shu, et al., ``TENET: An
Efficient Sparsity-Aware LUT-Centric Architecture for Ternary LLM Inference on
Edge,'' arXiv:2509.13765, 2025.

\bibitem{posit_accel2024} N.~Nakasato, Y.~Murakami, F.~Kono, and M.~Nakata,
``Evaluation of POSIT Arithmetic with Accelerators,'' arXiv:2401.14117, 2024.

\bibitem{pnnl_lns2025} ``Bridging the Gap Between LLMs and LNS with Dynamic Data
Format and Architecture Codesign,'' Pacific Northwest National Laboratory,
January 2025. No arXiv or DOI identifier was located; cited from the PNNL
publication record.

\bibitem{akhtaruzzaman2024phi} M.~Akhtaruzzaman, J.~U.~Tanvin, A.~A.~Shafie,
F.~Shahryer, and S.~Halder, ``Application of Phi ($\Phi$), the Golden Ratio, in
Computing: A Systematic Review,'' \emph{IEEE Access}, vol.~13, 2024.

\bibitem{schelpe2026grafting} S.~Schelpe, ``Smarter and Cheaper at Once: Byte-Exact KV-Cache Grafting Turns a Frozen Small Model into a Verified-Knowledge Flywheel,'' arXiv:2607.14431, 15~July 2026. \url{https://arxiv.org/abs/2607.14431}

\bibitem{zhu2026heal} Z.~Zhu, L.~Thai, S.~Yu, Y.~Liu, Y.~Qiao, C.~Wang, H.~Xu and J.~Shu, ``Demystifying Numerical Instability in LLM Inference: Achieving Reproducible Inference for Mission-Critical Tasks with HEAL,'' arXiv:2606.21023. \url{https://arxiv.org/abs/2606.21023}

\bibitem{tran2026rtlrepair} H.~T.~Tran, ``Open-Source LLM-Driven Formal Verification: A Multi-Agent Pipeline for RTL Repair,'' arXiv:2607.28877, 30~July 2026. \url{https://arxiv.org/abs/2607.28877}

\bibitem{cheng2026hint} T.~Cheng, Y.~Liu, D.~Zuo, Z.~Shi, H.~Zhang, X.~Hu,
M.~He, H.~Yan and Q.~Xu, ``HINT: Toward an Executable Hardware-Intent
Representation Layer for LLM-Driven RTL Generation,'' arXiv:2608.07625, 7~August
2026. \url{https://arxiv.org/abs/2608.07625}

\end{thebibliography}
